\documentclass[11pt]{article}

\usepackage[margin=1in]{geometry}
\usepackage{amsmath,amssymb,amsthm}
\usepackage{enumitem}
\usepackage{hyperref}
\usepackage{xcolor}
\usepackage{booktabs}

\hypersetup{colorlinks=true, linkcolor=blue, urlcolor=blue, citecolor=blue}

\title{TOAE209/TOAH209 -- Design and Analysis of Algorithms\\[4pt]
\large Week 8: Theoretical Questions}
\author{Foreign Trade University}
\date{}

\begin{document}
\maketitle

\section*{Instructions}

Part A contains 17 questions requiring written explanations or short proofs. Part B is a
10-question quiz (true/false and multiple choice). Attempt every question on your own
before consulting Part C (Solutions) at the end of this document.

\section*{Part A: Theory Questions}

\begin{enumerate}[label=\textbf{A\arabic*.}, leftmargin=2.2em]

\item Define \emph{optimal substructure} and \emph{overlapping subproblems}. Explain why
\emph{both} properties are needed for dynamic programming to be both correct and efficient,
and contrast this with divide-and-conquer (e.g.\ merge sort), which has optimal
substructure but \emph{not} overlapping subproblems.

\item Compare \emph{top-down memoization} and \emph{bottom-up tabulation}: how does each
avoid recomputing subproblems, what are the practical advantages of each, and why do they
have the same asymptotic running time?

\item Explain why the naive recursion for the Fibonacci numbers $F(n) = F(n-1) + F(n-2)$
makes exponentially many calls, while a memoized version makes only $O(n)$. Precisely how
many calls does the naive recursion make for $F(n)$ (in terms of Fibonacci numbers)?

\item State the Weighted Interval Scheduling problem precisely. Explain why the
earliest-finish-time greedy algorithm (optimal for \emph{unweighted} interval scheduling)
is \emph{not} optimal here, and why sorting by weight also fails.

\item Define the function $p(j)$ used in Weighted Interval Scheduling (after sorting by
finish time). Explain how to compute all values $p(1), \ldots, p(n)$ in $O(n \log n)$ time,
and why binary search is applicable.

\item Derive the Weighted Interval Scheduling recurrence $M[j] = \max(M[j-1],\, w_j +
M[p(j)])$ by reasoning about whether interval $j$ is in the optimal solution. Justify why
the ``take $j$'' branch may only use intervals $1, \ldots, p(j)$.

\item Prove that the recurrence in A6 correctly computes the maximum total weight of a
compatible subset of $\{1, \ldots, j\}$, by induction on $j$.

\item Describe how to reconstruct the actual optimal set of intervals (not just its value)
from the array $M$. State the exact condition under which interval $j$ is included during
the traceback, and the running time of the reconstruction.

\item State the 0/1 Knapsack problem. Write the 2-D DP recurrence for
$\mathit{dp}[i][w]$ (best value using the first $i$ items with capacity $w$), including the
base case and the two cases (item fits vs.\ does not fit), and explain the take-vs-skip
decision.

\item The 0/1 Knapsack DP runs in $O(nW)$ time. Explain carefully why this is
\emph{pseudo-polynomial} rather than polynomial. Your answer should refer to the number of
bits used to encode $W$ and explain how the running time behaves when that bit-length
grows.

\item State the Subset Sum problem and give its Boolean DP recurrence. Explain why, in the
space-optimized 1-D version, the inner loop over the target $t$ must run from high to low
(rather than low to high), and what would change if you reversed it.

\item Describe how to reconstruct which items were chosen in the 0/1 Knapsack, given the
filled 2-D table $\mathit{dp}$. State the exact traceback condition for ``item $i$ was
taken.''

\item Contrast the 0/1 knapsack with the \emph{unbounded} knapsack (each item reusable). In
the 1-D space-optimized implementations, exactly what differs between the two (loop
direction / order), and why does that single difference enforce ``at most once'' vs.\
``any number of times''?

\item Give the DP recurrences for both (a) the \emph{minimum number of coins} to make an
amount $A$ and (b) the \emph{number of distinct combinations} making $A$. For the counting
version, explain why looping over coins on the \emph{outside} (and amounts inside) counts
combinations rather than permutations.

\item State the Rod Cutting problem and its DP recurrence. Explain why it is an instance of
\emph{unbounded} knapsack, and describe how to recover the list of cut lengths (not just
the maximum revenue).

\item State the Longest Increasing Subsequence (LIS) problem and give the $O(n^2)$ DP
recurrence for $L[i]$, the length of the longest increasing subsequence \emph{ending at}
index $i$. Why is the final answer $\max_i L[i]$ rather than $L[n]$?

\item \emph{Discussion (open-ended).} You are given a new optimization problem and suspect
it has a DP solution. Describe a systematic process for designing and validating the DP:
how you would identify the subproblem and decision, write and prove the recurrence, choose
top-down vs.\ bottom-up, analyze the running time (and check whether it is genuinely
polynomial or only pseudo-polynomial), and empirically verify correctness against brute
force.

\end{enumerate}

\newpage
\section*{Part B: Quiz}

\begin{enumerate}[label=\textbf{B\arabic*.}, leftmargin=2.2em]

\item \textbf{(True/False)} Dynamic programming requires that a naive recursive solution
would re-solve the same subproblems many times (overlapping subproblems); otherwise plain
divide-and-conquer is already efficient.

\item \textbf{(Multiple choice)} In Weighted Interval Scheduling (sorted by finish time),
$p(j)$ is:
\begin{enumerate}[label=(\alph*)]
    \item the latest interval $i < j$ that overlaps $j$
    \item the latest interval $i < j$ that is compatible with $j$ (or $0$ if none)
    \item the interval with the largest weight among $1, \ldots, j-1$
    \item the earliest interval compatible with $j$
\end{enumerate}

\item \textbf{(Short answer)} Write the Weighted Interval Scheduling recurrence for $M[j]$
in terms of $M[j-1]$, $w_j$, and $M[p(j)]$.

\item \textbf{(True/False)} The 0/1 Knapsack DP running time $O(nW)$ is a polynomial-time
algorithm in the size of the input.

\item \textbf{(Multiple choice)} In the space-optimized 1-D Subset Sum / 0/1 Knapsack DP,
the inner loop over the capacity (or target) is iterated from high to low in order to:
\begin{enumerate}[label=(\alph*)]
    \item improve cache performance
    \item ensure each item is used at most once
    \item allow each item to be used unlimited times
    \item avoid integer overflow
\end{enumerate}

\item \textbf{(True/False)} Subset Sum is the special case of 0/1 Knapsack in which each
item's value equals its weight and we only ask whether the capacity can be filled exactly.

\item \textbf{(Short answer)} For the unbounded coin-change \emph{counting} DP, should the
coin loop be on the outside or inside of the amount loop in order to count distinct
combinations (order-independent)?

\item \textbf{(Multiple choice)} Rod Cutting is best classified as:
\begin{enumerate}[label=(\alph*)]
    \item 0/1 knapsack (each length usable at most once)
    \item unbounded knapsack (each length usable any number of times)
    \item weighted interval scheduling
    \item a greedy problem with no DP needed
\end{enumerate}

\item \textbf{(Short answer)} In the $O(n^2)$ LIS DP, $L[i]$ is defined as the length of
the longest increasing subsequence ending at index $i$. What is the recurrence for $L[i]$?

\item \textbf{(True/False)} Top-down memoization and bottom-up tabulation, applied to the
same recurrence, always compute the same optimal value and have the same asymptotic running
time, although memoization may skip subproblems that are never reached.

\end{enumerate}

\newpage
\section*{Part C: Solutions}

\subsection*{Part A}

\begin{enumerate}[label=\textbf{A\arabic*.}, leftmargin=2.2em]

\item \textbf{Optimal substructure} means an optimal solution to the whole problem is built
from optimal solutions to subproblems: once we fix one decision, what remains is a smaller,
independent instance of the same problem whose optimum is used by the global optimum. This
is what makes the \emph{recurrence correct}. \textbf{Overlapping subproblems} means the
naive recursion would solve the same subproblem repeatedly; storing each answer once is
what makes the recurrence \emph{efficient}. Both are needed: without optimal substructure
the recurrence would be wrong; without overlapping subproblems, memoization saves nothing.
Merge sort has optimal substructure (sorted halves combine into a sorted whole) but its
subproblems are \emph{disjoint} (different array ranges), so there is nothing to reuse --
plain divide-and-conquer already runs in $O(n \log n)$ and DP offers no advantage.

\item \textbf{Top-down memoization} writes the recurrence as a recursive function and
caches each subproblem's result on first computation, returning the cache on later calls;
subproblems are solved \emph{lazily} (only those reached). \textbf{Bottom-up tabulation}
fills a table iteratively in an order where each entry depends only on already-computed
ones. Memoization is usually easier to write (it mirrors the recurrence) and avoids
unreachable subproblems; tabulation has no recursion overhead, makes time/space transparent,
and often enables space optimizations (keeping only the last row). They have the same
asymptotic running time because both compute each distinct subproblem exactly once, and the
total time is (number of distinct subproblems) $\times$ (work per subproblem) in either
case.

\item Naive $F(n)$ recursion re-solves overlapping subproblems: computing $F(n)$ calls
$F(n-1)$ and $F(n-2)$, and $F(n-1)$ again calls $F(n-2)$, so $F(n-2)$ is computed twice,
$F(n-3)$ three times, and so on -- the call tree branches like the Fibonacci numbers
themselves and has exponentially many leaves. Memoizing computes each $F(k)$ for $k =
0,\ldots,n$ exactly once, so $O(n)$ work. The exact number of calls the naive recursion
makes for $F(n)$ is $2F(n+1) - 1$ (this is verified empirically in Problem 12).

\item \textbf{Input:} $n$ requests, request $i$ with start $s_i$, finish $f_i$ ($s_i <
f_i$), and weight $w_i \ge 0$; two requests are compatible iff their half-open intervals do
not overlap. \textbf{Objective:} a compatible set of maximum total weight. The
earliest-finish-time greedy is optimal for the \emph{count} (unweighted case) but ignores
weights: a single high-value interval can outweigh several low-value compatible ones it
conflicts with, so maximizing count is not the same as maximizing weight. Sorting by weight
fails too -- a heavy interval may block two even-heavier compatible intervals -- so no
simple greedy ordering is correct, and we use DP.

\item After sorting requests by finish time ($f_1 \le \cdots \le f_n$), $p(j)$ is the
largest index $i < j$ with $f_i \le s_j$ (the latest interval that finishes before $j$
starts), or $0$ if none exists. Since finish times are sorted, the set $\{ i : f_i \le s_j
\}$ is a prefix of the index range, so its largest element can be found by \textbf{binary
search} for $s_j$ in the sorted finish-time array. Each search is $O(\log n)$ and there are
$n$ of them, so all $p(j)$ are computed in $O(n \log n)$ (the same as the sort). This is
Problem 1.

\item Consider whether interval $j$ (the latest-finishing of $\{1,\ldots,j\}$) is in an
optimal solution $O_j$. If $j \in O_j$, then no other member of $O_j$ overlaps $j$; since
intervals are finish-sorted, every other member has index $\le p(j)$ (the latest interval
compatible with $j$), so $O_j \setminus \{j\}$ is a compatible subset of $\{1,\ldots,p(j)\}$
-- giving value $w_j + M[p(j)]$. If $j \notin O_j$, then $O_j$ is a compatible subset of
$\{1,\ldots,j-1\}$, giving value $M[j-1]$. The optimum is the larger: $M[j] = \max(M[j-1],
\, w_j + M[p(j)])$. The ``take'' branch may only use $1,\ldots,p(j)$ precisely because any
interval with index in $(p(j), j)$ overlaps $j$ and so cannot coexist with it.

\item By strong induction on $j$. \emph{Base:} $M[0] = 0$ is correct for the empty
instance. \emph{Step ($j \ge 1$):} let $O_j$ be an optimal compatible subset of
$\{1,\ldots,j\}$ with value $\mathrm{OPT}(j)$. If $j \notin O_j$, then $O_j$ is feasible for
$\{1,\ldots,j-1\}$, so $\mathrm{OPT}(j) \le M[j-1]$ by the inductive hypothesis; and any
optimum of $\{1,\ldots,j-1\}$ is feasible for $\{1,\ldots,j\}$, so $\mathrm{OPT}(j) \ge
M[j-1]$; thus this branch equals $M[j-1]$. If $j \in O_j$, then (as in A6) $O_j \setminus
\{j\}$ is feasible for $\{1,\ldots,p(j)\}$ with value $\mathrm{OPT}(j) - w_j \le M[p(j)]$,
so $\mathrm{OPT}(j) \le w_j + M[p(j)]$; and taking an optimum of $\{1,\ldots,p(j)\}$ plus
$j$ achieves $w_j + M[p(j)]$, so this branch is attainable. Hence $\mathrm{OPT}(j) =
\max(M[j-1],\, w_j + M[p(j)]) = M[j]$.

\item Trace back from $j = n$. Interval $j$ is in the optimum iff the ``take'' branch was at
least as good as the ``skip'' branch, i.e.\ iff $w_j + M[p(j)] \ge M[j-1]$. When this
holds, output $j$ and jump to $j \gets p(j)$ (the latest compatible predecessor); otherwise
move to $j \gets j-1$. Repeat until $j = 0$. Each step decreases $j$, so the reconstruction
is $O(n)$. (Ties may be broken either way; both yield optimal sets.) This is Problem 3.

\item \textbf{Input:} $n$ items with integer weights $w_i \ge 0$ and values $v_i \ge 0$,
and integer capacity $W$; choose a subset of maximum total value with total weight $\le W$,
each item at most once. \textbf{Recurrence:} $\mathit{dp}[i][w]$ = best value using items
$1,\ldots,i$ with capacity $w$:
\[
\mathit{dp}[0][w] = 0; \quad
\mathit{dp}[i][w] = \mathit{dp}[i-1][w] \ \text{if } w_i > w; \quad
\mathit{dp}[i][w] = \max(\mathit{dp}[i-1][w],\, v_i + \mathit{dp}[i-1][w - w_i]) \
\text{otherwise.}
\]
If item $i$ does not fit ($w_i > w$) we must skip it. Otherwise we choose the better of
skipping it (keep $\mathit{dp}[i-1][w]$) or taking it (gain $v_i$ and recurse on the
reduced capacity $w - w_i$ over the first $i-1$ items).

\item The size of the input -- the number of bits to encode it -- is about $O(n \log W)$,
since $W$ is written in binary using $\approx \log_2 W$ bits (and likewise for each $w_i,
v_i$). The running time $O(nW)$ is polynomial in the \emph{numeric value} $W$, but $W =
2^{\log_2 W}$ is \emph{exponential} in its bit-length. So if we increase the bit-length of
$W$ by one (doubling $W$), the running time doubles; increasing the bit-length by a factor
of two squares $W$ and thus squares the running time. An algorithm polynomial in the
numeric values but exponential in the input's bit-length is \textbf{pseudo-polynomial}. It
is efficient only when $W$ is small (polynomially bounded in $n$); 0/1 Knapsack is NP-hard,
so no truly polynomial algorithm is known.

\item \textbf{Subset Sum:} given non-negative integers $a_1,\ldots,a_n$ and target $T$,
decide whether some subset sums to exactly $T$. \textbf{Recurrence:} $R[i][t]$ = ``some
subset of the first $i$ numbers sums to $t$,'' with $R[0][0]=\text{true}$, $R[0][t>0] =
\text{false}$, and $R[i][t] = R[i-1][t] \lor (t \ge a_i \land R[i-1][t-a_i])$. In the 1-D
version with a single array $\texttt{reachable}[\,]$, applying $a_i$ by iterating $t$ from
$T$ \emph{down to} $a_i$ ensures that when we read $\texttt{reachable}[t - a_i]$ it still
reflects the state \emph{before} $a_i$ was applied -- so $a_i$ contributes at most once. If
we iterated low-to-high, $\texttt{reachable}[t-a_i]$ might already include $a_i$, letting
$a_i$ be used multiple times -- which would solve the \emph{unbounded} version instead.

\item Trace back from $\mathit{dp}[n][W]$ with $w = W$. For $i = n$ down to $1$: item $i$
was taken iff $\mathit{dp}[i][w] \ne \mathit{dp}[i-1][w]$ (taking it strictly changed the
value, meaning the ``take'' branch won). If taken, record $i$ and set $w \gets w - w_i$;
otherwise leave $w$ unchanged. Continue to row $i-1$. This is $O(n)$ and is Problem 7.

\item In \emph{0/1} knapsack each item is used at most once; in \emph{unbounded} knapsack
each may be used any number of times. In the 1-D space-optimized implementations the
\emph{only} difference is the inner loop direction (or, for counting, the loop nesting):
0/1 iterates the capacity/target from \emph{high to low}, so an item's update reads the
table state from \emph{before} this item was applied (each item contributes at most once);
unbounded iterates from \emph{low to high}, so an item's update may read entries already
updated by the same item in this pass (each item may contribute repeatedly).

\item (a) \textbf{Minimum coins:} $\mathit{dp}[0]=0$ and $\mathit{dp}[a] = 1 + \min_{c_i
\le a} \mathit{dp}[a - c_i]$ (or $+\infty$ if no coin fits), answer $\mathit{dp}[A]$ (or
$-1$ if $+\infty$). (b) \textbf{Counting combinations:} $\mathit{ways}[0]=1$ and, looping
over coins on the outside, $\mathit{ways}[a] \mathrel{+}= \mathit{ways}[a - c]$ for $a = c
\ldots A$. Putting the coin loop \emph{outside} commits coins in a fixed order, so each
combination (multiset) is generated exactly once; if amounts were the outer loop, the same
multiset would be counted in every order, counting \emph{permutations} (sequences) instead
of combinations.

\item \textbf{Rod Cutting:} given a rod of length $n$ and prices $p_\ell$ for pieces of
length $\ell$, cut into integer pieces to maximize total revenue. \textbf{Recurrence:}
$\mathit{rev}[0]=0$, $\mathit{rev}[L] = \max_{1 \le \ell \le L}(p_\ell + \mathit{rev}[L -
\ell])$. It is unbounded knapsack: ``item $\ell$'' has weight $\ell$, value $p_\ell$,
capacity $n$, and each length may be cut any number of times. To recover the cuts, store
for each $L$ the cut length $\ell^*(L)$ achieving the maximum; then starting from $L = n$,
repeatedly output $\ell^*(L)$ and set $L \gets L - \ell^*(L)$ until $L = 0$ (Problem 10).

\item \textbf{LIS:} given $a_1,\ldots,a_n$, find the length of the longest strictly
increasing subsequence (order preserved, elements may be dropped). \textbf{Recurrence:}
$L[i] = 1 + \max(\{L[j] : j < i,\, a_j < a_i\} \cup \{0\})$, i.e.\ extend the best
compatible earlier subsequence, or start fresh with length 1. The answer is $\max_i L[i]$,
not $L[n]$, because the optimal subsequence may \emph{end} at any index, not necessarily
the last one (e.g.\ if the largest values appear in the middle of the sequence). This is
Problem 11.

\item A systematic process: (1) \textbf{Identify the subproblem and decision.} Ask ``what
is the last/first choice the optimal solution makes, and what independent smaller instance
remains after it?'' Name an array indexed by the relevant parameters. (2) \textbf{Write and
prove the recurrence.} Express the optimum as a max/min over the possible decisions, then
prove correctness by induction (an optimal solution restricted to the subproblem is optimal
for the subproblem -- optimal substructure). (3) \textbf{Choose top-down vs.\ bottom-up.}
Memoize if the recurrence is more natural recursively or many subproblems are unreachable;
tabulate if you want transparent time/space or a space optimization. (4) \textbf{Analyze
the running time} as (number of subproblems) $\times$ (work per subproblem), and check
whether it is genuinely polynomial in the \emph{input size} or only pseudo-polynomial
(polynomial in a numeric value like $W$ or $A$). (5) \textbf{Verify empirically.} Implement
a brute-force solver (exponential, for small inputs) and compare against the DP on many
random instances and edge cases -- exactly as Problems 5, 6, 8, 9, 10, and 11 do this week.
A single mismatch should be minimized by hand to expose the flaw in the recurrence.

\end{enumerate}

\subsection*{Part B}

\begin{enumerate}[label=\textbf{B\arabic*.}, leftmargin=2.2em]

\item \textbf{True.} Without overlapping subproblems there is nothing to reuse, and plain
divide-and-conquer is already efficient (e.g.\ merge sort). DP's advantage comes precisely
from caching answers to subproblems that recur.

\item \textbf{(b)} the latest interval $i < j$ that is compatible with $j$ (finishes at or
before $j$ starts), or $0$ if none. This is found by binary search on the sorted finish
times.

\item $M[j] = \max\bigl(M[j-1],\; w_j + M[p(j)]\bigr)$, with $M[0] = 0$.

\item \textbf{False.} $O(nW)$ is \emph{pseudo}-polynomial: it is polynomial in the numeric
value $W$, but $W$ is exponential in its bit-length ($\approx \log_2 W$ bits), so the
algorithm is exponential in the input size. 0/1 Knapsack is NP-hard.

\item \textbf{(b)} to ensure each item is used at most once. Iterating high-to-low means
$\mathit{dp}[w - w_i]$ still holds the pre-item value, so the item contributes once;
low-to-high would allow reuse (the unbounded version).

\item \textbf{True.} Subset Sum is exactly 0/1 Knapsack with $v_i = w_i$ and $W = T$,
asking only whether the optimum equals (can reach) $T$.

\item \textbf{Outside.} The coin loop must be on the outside (amounts inside) to count
order-independent combinations; swapping the order would count ordered sequences
(permutations).

\item \textbf{(b)} unbounded knapsack -- each piece length may be used any number of times,
with capacity equal to the rod length.

\item $L[i] = 1 + \max\bigl(\{L[j] : j < i \text{ and } a_j < a_i\} \cup \{0\}\bigr)$; the
answer is $\max_i L[i]$.

\item \textbf{True.} Both evaluate the same recurrence, computing each distinct subproblem
once, so they return the same value with the same asymptotic running time; memoization
additionally skips subproblems that are never reached from the top-level call.

\end{enumerate}

\end{document}
